%── 5. Hydrogen in the power sector: from electrified heat to the merit order ──
\section{Hydrogen in the Power Sector: From Electrified Heat to the Merit Order}
\label{sec:power}

The heating merit-order analysis of Section~\ref{sec:results:meritorder} located hydrogen inside
the heating stack; on operating cost, hydrogen undercuts the carbon-priced fossil competitor at
the building and district-heat peak in the high-carbon scenarios, while on base-load levelised
cost it leads the heat pump in the cold, high-electricity-price north-west and trails it across
the warm south and east. This section asks a distinct question about the power system that carries the
electrified heat. As heating electrifies, a rising share of heat demand becomes electricity demand,
and the decarbonised grid that serves it still requires firm capacity for the cold-snap peak. We
therefore ask, on the same operating-cost basis a system operator actually uses, where hydrogen
sits in the \emph{power} merit order, and how large a role that position implies.

This question is the buildings-to-power dimension of the third research question (RQ3),
the synergies between hydrogen in buildings and the wider energy system, and it is the
one cross-sector coupling this paper takes up. The scope is deliberately bounded to that
coupling. The further synergies with industry and transport are out of scope and are
deferred to subsequent work, as set out in Section~\ref{sec:rq}. Throughout, electricity
is priced as wholesale markets price it, at the short-run marginal cost of the marginal
generator, because capital is a sunk cost that does not set the spot price; the
investment question is returned to, as a caveat, in Section~\ref{sec:power:caveat}.

The dispatchable layer is resolved by a reduced-form merit order, which sets the fleet
size and the running hours quoted throughout, and is cross-checked against an explicit
unit-commitment optimisation of the firm dispatchable fleet (firm CCGT against
open-cycle gas or hydrogen peakers, with minimum stable output, ramp and minimum up and
down times, a 5~per cent operating reserve, start-up costs and storage state of charge;
Appendix~\ref{sec:method:merit}). With the commitment constraints disabled the two agree
to within 0.33~per cent on dispatched energy and 0.40~per cent on peak capacity in 28 of
the 29 countries; Romania is the exception, and
Appendix~\ref{sec:results:dispatch} sets out why. The exercise tests where
the heating result lands in the power stack and bounds the implied role.

Electrified heating represents a material but bounded share of 2050 EU power demand. Across all 29
markets, heating electricity is about 9~per cent of total
power demand (Figure~\ref{fig:electrified}b), large enough to shape the load but not to dominate it.
This is the scaled load that determines where hydrogen competes. The cold-snap peak of that demand
block is the segment the power system must cover with firm dispatchable capacity, and it is there
that hydrogen enters the merit order.

The component shares behind that 9~per cent follow from the degree of electrification. By 2050
heating is heavily electrified. Heat pumps and resistance heating meet about 41~per cent of useful
heat directly under H2 Push, rising to about 51~per cent once the electric portion of district heat
(large heat pumps), taken as its share of total useful heat, without being added in full, is included,
and about 51~per cent directly under Net Zero (Figure~\ref{fig:electrified}a). Converting those
shares to electricity through the temperature-dependent coefficient of performance and adding the
result to a national baseline load grown to 2050 yields the 9~per cent share of total power demand
above.

\begin{figure}[H]
    \centering
    \caption{The bridge from electrified heating to 2050 electricity demand, under H2~Push,
    all 29 markets. The red dashed line in each panel is the study-area average. Panel~(a)
    counts district heat at its 0.45 electrification factor, which is why its 51~per~cent
    exceeds the 41~per~cent direct-electrification average of Table~\ref{tab:electrified}.}
    \includegraphics[width=0.49\textwidth]{paper/P28_electrified_heat.pdf}
    \hfill
    \includegraphics[width=0.49\textwidth]{paper/P29_heat_electricity_share.pdf}
    \label{fig:electrified}
    \\[2pt]{\footnotesize\textit{Source: Authors' contribution.}}
\end{figure}

The same picture, country by country and across the decade, is set out in
Table~\ref{tab:electrified}. The heat-weighted EU-27+UK+CH share of heat met directly by electricity
climbs from about 23~per cent in 2030 to about 31~per cent in 2040 and about 41~per cent in 2050,
and once the electric share of district heat is added the 2050 figure rises to about 51~per cent.

\begin{table}[htbp]
\centering
\small
\caption{Share of residential heat met directly by electricity (heat pumps and resistance heating), per country, 2030/2040/2050, under H2 Push, with the heat-weighted EU-27+UK+CH average in the final row.}
\label{tab:electrified}
\begin{tabular}{lrrr}
\toprule
Country & 2030 & 2040 & 2050 \\
\midrule
NL & 24 & 29 & 46 \\
FI & 45 & 46 & 44 \\
DK & 27 & 35 & 44 \\
SE & 41 & 44 & 43 \\
UK & 27 & 29 & 43 \\
PL & 10 & 31 & 42 \\
CH & 30 & 34 & 42 \\
IE & 25 & 28 & 41 \\
HU & 10 & 29 & 41 \\
LT & 19 & 34 & 41 \\
LV & 12 & 33 & 41 \\
CZ & 15 & 32 & 41 \\
EE & 20 & 36 & 41 \\
SI & 20 & 33 & 40 \\
RO & 10 & 31 & 40 \\
FR & 36 & 36 & 40 \\
DE & 23 & 28 & 40 \\
AT & 27 & 32 & 40 \\
BE & 11 & 27 & 39 \\
HR & 11 & 30 & 39 \\
IT & 12 & 28 & 39 \\
SK & 15 & 30 & 39 \\
LU & 23 & 27 & 38 \\
BG & 23 & 37 & 38 \\
ES & 24 & 32 & 37 \\
PT & 42 & 41 & 37 \\
EL & 18 & 28 & 35 \\
MT & 79 & 52 & 32 \\
CY & 37 & 31 & 29 \\
\midrule
EU-27+UK+CH (heat-weighted) & 23 & 31 & 41 \\
\bottomrule
\end{tabular}
\\[2pt]{\footnotesize\textit{Source: Authors' contribution, scenario Monte Carlo medians.}}
\end{table}


\subsection{Where Hydrogen Sits in the Power Merit Order}
\label{sec:power:meritorder}

Hydrogen competes for the peaker role on short-run marginal cost, the fuel-plus-carbon
operating cost that a system operator dispatches on, with no capital charge. Under
H2~Push, which pairs the high carbon path with the rapid hydrogen cost decline, hydrogen
undercuts the gas peaker on operating cost in \textbf{seven of 29 countries}
(Figures~\ref{fig:meritladder} and~\ref{fig:crossover}), exactly the salt-cavern set, a
point we return to below because the winning set being identical to the classification
means the seven is conditional on it rather than independent of it,
implying it is dispatched ahead of gas at the cold-snap peak in those markets. That
count holds both turbines at the same 0.40 efficiency, which is the even-handed
comparison and the one we carry throughout.

On the \emph{evolving-efficiency} series, in which the gas turbine reaches 0.42 by 2050
and the hydrogen turbine 0.48 under H2~Push, so the hydrogen unit is deliberately given
the better machine, the count rises to 15. We report both because they answer different
questions, the first what a plausible 2050 fleet would dispatch and the second whether
hydrogen wins on fuel cost alone; the seven that survive the even-handed test are the
same seven that win decisively on the looser one. This finding is consistent with the
system-value literature on complementary firm and flexible low-carbon resources
\citep{Bistline2020ValueOfTechnology,Lotze2024Complementary}, which locates hydrogen's
contribution at the dispatchable margin, and not in bulk energy. The count is
contingent on the 2050 carbon-price path, which is the dominant lever on the ranking
(Appendix~\ref{sec:sensitivity}). A lower carbon price narrows the gas peaker's
operating-cost penalty and would reduce the count, a higher one would raise it.

The advantage is, however, sharply bimodal and concentrated where seasonal storage is cheap. In the
seven salt-cavern markets, the United Kingdom, Denmark, the Netherlands, Romania, Poland, Germany
and France, hydrogen undercuts the gas peaker by a decisive \euro{}85 to 99/MWh-e on this evolving-efficiency series, and by \euro{}72 to 89/MWh-e on the equal-efficiency series the headline count uses; in the remaining eight of the 15 the margin is only \euro{}0 to 8/MWh-e, within the noise of the cost estimate.
The operative result is therefore better read as about seven markets with a material operating-cost
advantage, all of them cavern-endowed, and a further handful at the margin. It is not a uniform
15 (Figures~\ref{fig:crossover} and~\ref{fig:cavern7}).

The mechanism is the comparison of two operating costs. At the cold-snap peak the price is set by
the dearest dispatched generator, and the candidate firm units rank by short-run marginal cost:
\begin{equation}
    c^{\text{gas}}_{\text{SRMC}} = \frac{p^{\text{gas}} + \varphi\, e_g}{\eta_g},
    \qquad
    c^{\text{H}_2}_{\text{SRMC}} = \frac{p^{\text{H}_2}}{\eta_h} + \sigma,
    \label{eq:srmc}
\end{equation}
where each short-run marginal cost is expressed in EUR/MWh-e of electricity sent out:
$p^{\text{gas}}$ and $p^{\text{H}_2}$ are the delivered fuel prices (EUR/MWh of fuel),
$\varphi$ is the carbon price (EUR/tCO$_2$), $e_g$ the gas emission factor (tCO$_2$/MWh
of fuel) and $\eta_g,\eta_h$ the dimensionless electric efficiencies of the gas and
hydrogen turbines. The first term of each expression is the fuel-and-carbon operating
cost; both exclude capital, which is sunk and does not set the spot price (the capital
recovery of the hydrogen unit is the separate test of Section~\ref{sec:power:caveat}).
The seasonal-storage adder $\sigma$ is the annualised cost of holding and cycling
hydrogen in seasonal store; it is purely a storage-and-cycling charge and excludes
transmission losses. It is approximately \euro{}50/MWh-heat where caverns exist and approximately
\euro{}100/MWh-heat where they do not, equivalently \euro{}45 and \euro{}90 per MWh of
hydrogen at the 0.9 boiler efficiency used throughout, at rates consistent in order with
the bulk-storage literature \citep{Lord2014GeologicStorage,Papadias2021BulkStorage,Talukdar2024UHSEurope} and
applied on the cavern geology of \citet{Caglayan2020caverns}. Equation~\eqref{eq:srmc} needs it per MWh of
electricity sent out, so the implementation divides by the turbine efficiency, giving
about \euro{}112 and \euro{}225/MWh-e at $\eta_h=0.40$; it is the largest single driver of the hydrogen unit's cost gap, and its derivation is
set out under seasonal storage in Appendix~\ref{sec:method:storageadder}.

Ordered by Equation~\eqref{eq:srmc} against the low-carbon base of the 2050 fleet, the
merit order runs from near-zero variable renewables, through nuclear (about
\euro{}12/MWh) and biomass (about \euro{}55/MWh), up to the dispatchable tip. There the
hydrogen turbine is bimodal across countries, so a single average describes it poorly.
On the evolving-efficiency series used for capital recovery, which gives the hydrogen
turbine the better machine, at the cross-country median it sits at about
\euro{}238/MWh-e against a carbon-priced gas peaker at about \euro{}239/MWh-e, which is
a tie. Held at equal efficiency, the series that sets the win counts, the same medians
are about \euro{}286 against \euro{}256/MWh-e, so hydrogen sits clearly above gas in the
typical market. The cross-country mean of about \euro{}221/MWh-e, some \euro{}18/MWh-e
below gas, is pulled down by the cavern tail and should not be read as the typical
market. On that same evolving-efficiency series, in the seven cavern markets the
hydrogen turbine falls to about \euro{}141 to 155/MWh-e, decisively below gas; held at
equal efficiency the equivalent range is \euro{}168 to 184. Eight of the remaining 22
sit within \euro{}8/MWh-e of it, inside the noise of the estimate, and the other 14 sit
above it (Figure~\ref{fig:meritladder}).

\begin{figure}[H]
    \centering
    \caption{\textbf{2050 power merit order in Denmark, under H2 Push.} A salt-cavern
    market. Short-run marginal cost by technology, on the \emph{equal-efficiency} turbine
    series that holds both machines at the same conversion efficiency. That series prices the gas peaker at
    \euro{}256/MWh-e and the hydrogen turbine at \euro{}168 here; the evolving-efficiency
    series puts these same two Danish machines at \euro{}241 and \euro{}142, because it moves
    the gas machine's efficiency and fuel price as well as the hydrogen turbine's. The
    \euro{}239 and \euro{}221 quoted in the paragraph above are cross-country statistics,
    the median gas cost and the mean hydrogen cost, and are not Denmark's.
    Costs are Denmark's; the capacity column is
    the study-area 2050 fleet, and the two dispatchable rows carry the same firm fleet
    because they are alternative fuels for it. Across the 22 non-cavern markets the
    hydrogen turbine sits about \euro{}36/MWh-e above the gas peaker, so the median
    country places hydrogen above gas.}
    \includegraphics[width=0.88\textwidth]{paper/P26_merit_order.pdf}
    \label{fig:meritladder}
    \\[2pt]{\footnotesize\textit{Source: Authors' contribution.}}
\end{figure}

\begin{figure}[H]
    \centering
    \caption{\textbf{Hydrogen's operating-cost advantage over the gas peaker, 2050.} By country, under H2~Push, on short-run marginal cost with capital excluded, on the \emph{evolving-efficiency} turbine series. Hydrogen clears zero in 15 markets, but the shading separates two cases the count runs together. Seven
    markets, all salt-cavern, lead by \euro{}85 to 99/MWh-e and carry the result reported
    in the text. Eight more clear zero by at most \euro{}8/MWh-e, which is inside any
    reasonable read of model noise, and they should not be counted with the seven.}
    \includegraphics[width=\textwidth]{paper/P27_h2_crossover.pdf}
    \label{fig:crossover}
    \\[2pt]{\footnotesize\textit{Source: Authors' contribution.}}
\end{figure}

\begin{figure}[H]
    \centering
    \caption{\textbf{Where hydrogen is decisively cheaper.} The seven salt-cavern
    markets, on the \emph{equal-efficiency} turbine series that holds both machines at the
    same conversion efficiency. This is the series the headline seven is reported on, so its margins,
    \euro{}72 to 89/MWh-e, are narrower than the \euro{}85 to 99/MWh-e that
    Figure~\ref{fig:crossover} shows for the same seven markets on the evolving-efficiency
    series, where the hydrogen turbine is credited with the better machine.}
    \includegraphics[width=0.78\textwidth]{paper/P33_cavern_markets.pdf}
    \label{fig:cavern7}
    \\[2pt]{\footnotesize\textit{Source: Authors' contribution.}}
\end{figure}

The 2050 snapshot is the end point of a trajectory. Table~\ref{tab:opcost} traces the two operating
costs from 2025. Hydrogen's falls steeply as delivered hydrogen cheapens and the turbine improves,
while the carbon-priced gas peaker rises, so the two cross only late and only under H2 Push.

\begin{table}[htbp]
\centering
\footnotesize
\caption{\textbf{Mean peaker operating cost, by scenario and year.} Across the EU-27, the United Kingdom and Switzerland; short-run marginal cost in EUR/MWh-e, fuel and carbon only, no capital. Hydrogen's cost falls as delivered hydrogen cheapens and the turbine improves; the carbon-priced gas peaker rises with the carbon path (about 70 to 350 EUR/tCO$_2$ in the two high-carbon scenarios). The two cross only late and only under H2 Push (221 versus 239 in 2050 on this MEAN basis; on the median the same year is a tie at 238 against 239, and the AE submission reports the median because this distribution is bimodal and the mean is dragged by the salt-cavern tail, so the sign of the crossing depends on which statistic is quoted); in the lower-carbon scenarios gas stays cheaper throughout. The EU-27+UK+CH mean understates the salt-cavern markets, where the hydrogen cost is far lower (Figure~\ref{fig:crossover}).}
\label{tab:opcost}
\begin{tabular}{llrrrr}
\toprule
Scenario & Peaker & 2025 & 2030 & 2040 & 2050 \\
\midrule
Current Policies & Gas & 114 & 121 & 141 & 143 \\
& H$_2$ & 667 & 605 & 487 & 416 \\
\addlinespace
Stated Policies & Gas & 127 & 147 & 165 & 167 \\
& H$_2$ & 667 & 417 & 316 & 281 \\
\addlinespace
Net Zero & Gas & 135 & 186 & 224 & 239 \\
& H$_2$ & 667 & 417 & 316 & 281 \\
\addlinespace
H2 Push & Gas & 135 & 186 & 224 & 239 \\
& H$_2$ & 667 & 393 & 267 & 221 \\
\addlinespace
\bottomrule
\end{tabular}
\\[8pt]{\footnotesize\textit{Source: Authors' contribution.}}
\end{table}


The same scenario monotonicity seen at the heating level recurs here.
Table~\ref{tab:arenas} collects the count of countries in which
hydrogen wins across the three arenas and the four scenarios. In the power merit order hydrogen wins nowhere
under Current or Stated Policies, but in seven countries under Net Zero and 15 under H2 Push,
precisely the high-carbon worlds. The same monotonicity appears at the building winter peak
(0, 2, 7, 14) and inside district heating (0, 0, 7, 7), both on the symmetric basis, and on
the asymmetric one at (0, 4, 9, 16) and (0, 5, 11, 20). The mechanism is the carbon pass-through
that raises the gas peaker's marginal cost \citep{Fabra2009etspassthrough}, sharpened by the
extension of carbon pricing to buildings and road transport under the second Emissions Trading
System, ETS2 \citep{EuropeanCommission2024ETS2, BNEF2025ETS2}; the position of hydrogen in the
stack is set as much by the fortunes of its fossil rival as by its own cost.

% [htbp], not [H]. With [H] this table was pinned directly under Table 4, the two
% together ran 87 pt past the bottom of the text block, and the note overprinted the page
% number. [H] cannot move a float that does not fit; it just lets it overrun. Allowing
% this one to float puts it at the top of the following page, which also fills the page
% that was left three-quarters empty behind it.
\begin{table}[htbp]
    \centering
    \caption{Operating-cost wins do not translate into capital self-financing. The
    power-sector row is on the evolving-efficiency turbine series; the even-handed series,
    holding both turbines at 0.40, gives seven wins under both Net Zero and H2~Push,
    exactly the salt-cavern markets.}
    \label{tab:arenas}
    \small
    \begin{tabular}{lrrrr}
        \toprule
        & Current & Stated & Net Zero & H2 Push \\
        \midrule
        \multicolumn{5}{l}{\textit{Panel A. Operating-cost wins (merit-order dispatch; count of 29)}} \\
        \addlinespace[2pt]
        \multicolumn{5}{l}{\emph{Symmetric basis: hydrogen pays its own last-mile adder, gas}} \\
        \multicolumn{5}{l}{\emph{enters district heat at the wholesale hub price. Headline.}} \\
        Building winter peak (vs gas peaker)                      &  0 & 2 &  7 & 14 \\
        District-heat peak (vs gas CHP)                           &  0 & 0 &  7 &  7 \\
        \addlinespace[3pt]
        \multicolumn{5}{l}{\emph{Asymmetric basis: hydrogen pays no last-mile network, gas is}} \\
        \multicolumn{5}{l}{\emph{priced at the residential retail tariff.}} \\
        Building winter peak                                      &  0 & 4 &  9 & 16 \\
        District-heat peak                                        &  0 & 5 & 11 & 20 \\
        \addlinespace[3pt]
        \multicolumn{5}{l}{\emph{Neither charge enters the power arena, so these serve both.}} \\
        Power peaker, evolving turbine efficiency                 &  0 & 0 &  7 & 15 \\
        Power peaker, both turbines held at 0.40                  &  0 & 0 &  7 &  7 \\
        \midrule
        \multicolumn{5}{l}{\textit{Panel B. Capital recovery on market rent (count of 29)}} \\
        Recovers peaker capital from market rent                  &  0 & 0 &  0 &  0 \\
        \bottomrule
    \end{tabular}

    % The note was centred prose inheriting the table's \centering, so every line
    % centred independently and the block sat wider than the table it annotates. Set it
    % ragged-right in a box the width of the text block, which is what every other note
    % in this paper does.
    \smallskip
    \begin{minipage}{\textwidth}
    \raggedright\footnotesize
    \textit{Cumulatively to 2050 the H2 Push hydrogen peaker loses about
    \euro{}357~bn against about \euro{}142~bn for an equivalent gas peaker, so neither fuel
    recovers its capital in any country (Section~\ref{sec:power:caveat}). District-heating wins
    in Panel A are countries in which hydrogen undercuts gas combined heat and power (CHP) on
    cost; gas CHP nonetheless remains the marginal source that sets the district-heat peak in 20
    of 29 countries (hydrogen in 9), because hydrogen is dispatched below it. Costs follow the
    Danish Energy Agency Technology Catalogue \citep{DEA2023technology} with the scenario carbon
    and hydrogen price paths (Section~\ref{sec:method}).}\par
    \smallskip
    \textit{Source: Authors' contribution.}\par
    \end{minipage}
\end{table}

The win counts of Table~\ref{tab:arenas} aggregate a heterogeneous country picture, and that
heterogeneity is the paper's first contribution. Figure~\ref{fig:winloss} resolves it,
mapping for each of the 29 countries whether hydrogen wins or loses in each of the three
arenas under H2 Push, and distinguishing the decisive salt-cavern wins from the marginal
non-cavern ones. The pattern is bimodal. The seven cavern-endowed markets (Germany, the
Netherlands, France, Denmark, the United Kingdom, Poland and Romania) win across all three
arenas and carry the firm-capacity tip, while the remaining wins are marginal and, as the
missing-money result will show, unfunded. So although hydrogen wins on operating cost in 14,
7 and 15 of 29 countries across the three arenas, only about seven of those wins are
decisive.

% [tbp], not [H]. This is a 29-row heat map that needs most of a page, and [H] pinned it
% directly under Table 5, where it ran off the bottom and over the folio. Excluding `h`
% keeps it at a page top, which is the only place it fits whole.
\begin{figure}[tbp]
    \centering
    \caption{Where hydrogen wins, by country and arena (H2 Push, 2050). The power-peaker
    column is the evolving-efficiency turbine series, which gives 15 wins; the
    equal-efficiency series that the headline count uses gives seven, all salt-cavern.
    Shading separates decisive wins, in the cavern markets, from marginal ones.}
    \includegraphics[width=0.80\textwidth]{paper/P35_winloss.pdf}
    \label{fig:winloss}
    \\[2pt]{\footnotesize\textit{Source: Authors' contribution.}}
\end{figure}

\subsection{The Role This Implies, a Capacity Role Measured in Peak Hours}
\label{sec:power:share}

How much of the power system can hydrogen take? The answer divides in two, and
the division \emph{is} the result. The firm peaking fleet, sized to a
loss-of-load standard in the hourly dispatch of Appendix~\ref{sec:results:dispatch}, is
about 278~GW across the EU-27, the United Kingdom and Switzerland, about 31~per cent of
the 889~GW system peak. Both figures, and every power-sector number that follows from
them, are computed on \emph{one realisation of a single synthetic weather year} and not
on reanalysis or metered weather. That is the largest simplification in this paper, it
is set out in full in Appendix~\ref{sec:limitations}, and a reader who wants a system
adequacy number should not take one from here. Four further qualifications belong with
those two figures, and two of them bound the 278~GW itself. The fleet is sized on a completely inflexible electrified
heat load. Making 10~per cent of the heat-pump load shiftable within three hours leaves
it at about 277~GW, and at the most generous corner we sweep, 30~per cent shiftable
across thirteen hours, three energy-conserving smoothing operators span 265.6 to
248.0~GW, a reduction of 4.3 to 10.7~per cent.

And the single synthetic cold-and-still event that sizes it is assumed to last six days;
sweeping that length from two days to ten moves the fleet sized on the inflexible load from 253.7 to
302.0~GW, so 278~GW carries a band of roughly $\pm$9~per cent from that assumption
alone, wider than the flexibility correction. Both make 278~GW, and the
\euro{}357~bn shortfall that scales with it, upper bounds. Neither is a point
estimate. The 889~GW sums 29
national peaks that do not coincide, so it bounds the requirement from above in an interconnected system, which is what the
copper-plate arm of Appendix~\ref{sec:results:dispatch} tests. The 278~GW is an
increment on a firm combined-cycle fleet of about 195~GW that the dispatch holds fixed
and runs ahead of the peaker, so total firm dispatchable capacity is about 473~GW, or
53~per cent of that peak; what follows concerns the peaking increment, which is the part
hydrogen competes for. Where hydrogen undercuts gas it can serve that backup tip;
aggregating the firm capacity of the 15 winning countries,
\begin{equation}
    \Phi_{\text{cap}}
    = \frac{\sum_{c\,:\,c^{\text{H}_2}_c < c^{\text{gas}}_c} K_c^{\text{pk}}}
           {\sum_c D_c^{\text{pk}}}
    \approx \frac{233~\text{GW}}{889~\text{GW}} \approx 26\%,
    \qquad
    \Phi_{\text{en}} = \frac{E^{\text{pk}}}{D^{\text{tot}}}
    \approx \frac{28~\text{TWh}}{4{,}513~\text{TWh}} \approx 0.63\%,
    \label{eq:share}
\end{equation}
where, summing over countries $c$, the subset $c:c^{\text{H}_2}_c < c^{\text{gas}}_c$ in
the numerator is the 15 countries in which hydrogen wins the dispatch on operating cost;
$K_c^{\text{pk}}$ is country $c$'s firm peaker capacity (GW), $D_c^{\text{pk}}$ its
annual peak electricity demand (GW), $E^{\text{pk}}$ the EU peaker energy (TWh) and
$D^{\text{tot}}$ total annual power demand (TWh).

One convention decides that 15, and we
have not settled it. The looser short-run cost series that carries the capital-recovery
accounting prices the molecule at 0.85 of its delivered price on the reasoning that a
turbine sited on the hydrogen backbone avoids the last-mile margin, while our own cost
layer treats it as though it does not. Removing the discount takes the count from 15 to 8, the seven cavern-endowed markets
plus Ireland, and with it $\Phi_{\text{cap}}$
from 26.3 to 23.3~per cent and $\Phi_{\text{en}}$ from 0.63 to 0.36~per cent. The 15 and
the shares above should therefore be read as the top of an 8 to 15 range on a convention
we do not resolve, and the seven-market equal-efficiency count that carries the headline
is unaffected, because it prices both turbines on the same basis.

The 4{,}513~TWh denominator rests on an assumption we state here.
Non-heat national electricity demand is set at its 2023 level, 2{,}926~TWh across the 29
markets on ENTSO-E and Ember reporting, and grown by a smooth factor reaching 1.40 by 2050
(1.00, 1.10 and 1.25 at 2025, 2030 and 2040) to stand for the electrification of road
transport and industry. That is an assumed path, equivalent to about 1.35~per cent a year compounding, which no model produces, and a single factor is applied across countries because
the comparison at issue is between two firm technologies inside a single demand envelope; demand futures are not being compared. Heat-pump electricity from the heating model is added on top,
at 417~TWh in 2050, giving the 4{,}513~TWh total. Both shares scale inversely with the
growth factor; a flat baseline would raise this winners' energy share to about 0.85~per cent, and the whole-fleet share to about 1.2~per cent and leave
the capacity share, set by the winter peak rather than by annual energy, materially
unchanged.

$\Phi_{\text{cap}}$ is thus the share of system peak capacity that hydrogen can serve and
$\Phi_{\text{en}}$ the share of generated energy it supplies. Hydrogen can take
about 26~per cent of study-area peak \emph{capacity} through the winning countries, which is some
84~per cent of the 278~GW peaker fleet, yet supply about 0.63~per cent of generated
\emph{energy}, because the peaker runs only about 139 full-load hours a year. Counting the whole fleet, and not the winners alone, gives the 31~per cent and 0.86~per cent quoted elsewhere in
this paper; the difference between the two pairs is the 14 non-winning markets, of
which 9 hold some peaker capacity. The two figures describe the same plant from the two sides of the system, capacity and energy. Hydrogen's electricity role is
accordingly a capacity role, and its energy contribution is negligible. It can serve as the clean firm backstop for the
electrified heat load while remaining a negligible share of generated kilowatt-hours. The reading
echoes the wider finding that the system value of a clean firm resource resides in the scarce hours it covers, which its annual output does not measure \citep{Bistline2020ValueOfTechnology, Quarton2019Value}.

\subsection{The Investment Caveat}
\label{sec:power:caveat}

Setting the price and taking the role on operating cost does not, on its own, finance the plant.
The missing-money test of Section~\ref{sec:results:missingmoney} applies unchanged to
the power peaker. The inframarginal rent recovers only about 9~per cent of the power
peaker's annualised capital by 2050, capacity-weighted across the five cavern markets
that build a hydrogen peaker, which average 80 running hours against 139 for the fleet as a whole
(Figure~\ref{fig:peakerbridge}).

Two recovery percentages appear. Figure~\ref{fig:peakerbridge} plots the study-area
aggregate, which reaches 14.6~per cent by 2050, and the 9~per cent quoted here weights
only the five cavern markets that build a hydrogen peaker, by the capacity each builds,
and Germany's 88~GW at 48 running hours dominates that weighting.

The aggregate is still rising at 2050. It turns positive in 2037 and climbs monotonically to
the horizon, so it bounds nothing beyond it and we make no claim about where it settles.

Cumulatively to 2050 the H2~Push hydrogen peaker loses \euro{}357.3~bn against
\euro{}142.4~bn for an equivalent gas peaker, on the reduced-form dispatch that also
sets the fleet size and the running hours. The unit-commitment layer runs the same
economics on its own dispatch and lands close to them, at losses of \euro{}358.4~bn and
\euro{}143.2~bn with its constraints switched off and \euro{}356.0~bn and
\euro{}142.3~bn with them on, because the constraints
raise peaker run-hours and the plant earns in more hours as well as burning more fuel.
That spread is the effect of the commitment constraints, reported in
Appendix~\ref{sec:results:dispatch}.

Because electricity
is priced on operating cost, this shortfall does not prevent hydrogen from setting the price or
winning the dispatch in the 15 countries; it means that a deliberate hydrogen-peaker build implies
a standing capacity payment or equivalent support, the classic missing-money outcome of low-utilisation peaking plant \citep{Joskow2008capacity,Newbery2016Missing}. The verdict is economic, not one about reliability. The modelled system remains adequate, with a loss of load averaging about two hours a
year (Appendix~\ref{sec:results:dispatch}), and the difficulty is paying for the firm capacity, not keeping the system adequate. The dominant driver of the gap is the seasonal hydrogen-storage adder
$\sigma$ in Equation~\eqref{eq:srmc}, which the delivered fuel price does not set; even at hydrogen's physical production floor the peaker's run-hours set the ceiling on
what market rent can recover, and its fuel cost does not.

Two qualifications belong with that verdict, and both cut against reading it as a
technology result. The first is that the shortfall is a property of the market design we impose,
and the machine itself does not cause it. Rent per kW is the scarcity price times the hours the plant runs, so with the scarcity
premium fixed at \euro{}60/MWh the recovery percentage follows from run-hours alone by
Eq.~\eqref{eq:powerrecovery} (Appendix~\ref{sec:method:merit}). Germany runs 48
hours, which at
\euro{}60/MWh is \euro{}2.9/kW-yr against \euro{}53.4/kW-yr of annualised capital, or
5.4~per~cent. The model also enforces a three-hour loss-of-load standard and prices
unserved energy at a value of lost load of \euro{}8{,}000/MWh, but pays that price to no
generator. Crediting the peaker for those three hours at the model's own value of lost
load adds about \euro{}24/kW-yr and lifts capacity-weighted recovery from 8.7 to about 51
per cent, with Denmark moving from 40 to about 87. The missing-money finding is therefore a statement about an energy-only market with an
administrative price cap, set out as such in Appendix~\ref{sec:limitations:energyonly}.
It does not say the plant cannot be paid for.

The second is that a capacity payment large enough to build firm capacity would not, on
these numbers, build \emph{hydrogen}. The hydrogen peaker's rent shortfall is a standing payment of \euro{}31 to 63/kW-yr across the five markets that build, against \euro{}41 to 52/kW-yr for a gas peaker in the same markets. The gas figure is the smaller of the two in four of those five, covering about 96~per cent of the winning capacity, because the hydrogen-ready machine costs about
\euro{}600/kW against about \euro{}450/kW and that capital gap outweighs its smaller rent
shortfall. A technology-neutral mechanism procures unabated gas, which is what the emissions limits in Article~22(4) of the Electricity Regulation are meant to prevent and, at 550~gCO$_2$/kWh and 350~kgCO$_2$/kWe-yr, do not \citep{EU2019ElectricityRegulation,ACER2023capacity}; the 2024 British four-year-ahead auction is a mechanism doing exactly this, procuring 27.3~GW of unabated gas out of 43.1~GW at \pounds{}60/kW-yr \citep{DESNZ2025T4auction}. Selecting hydrogen instead requires the payment to be conditioned on emissions, and the implied conditioning
premium is about \euro{}172/tCO$_2$ on top of the carbon price the gas plant already pays,
rising to roughly \euro{}239/tCO$_2$ if the gas peaker's negative rent is floored at zero
on the ground that no operator would hold that position. The direction is unchanged either
way, and the lower figure is the conservative end.

\begin{figure}[H]
    \centering
    \caption{The power peaker is competitive on operating cost yet unfunded, under H2~Push.
    Panel~(a) is the share of the hydrogen turbine's annualised capital its market rent
    returns, against the same ratio for a gas peaker, with the three other scenarios drawn
    faint for context. Its 15~per~cent is the study-area aggregate, the sum of rent over the sum of annualised
    capital across the countries the dispatch builds a peaker in. The 9~per~cent quoted in
    the text is a different statistic on the same run, the capacity-weighted mean of the
    per-country ratios across the five cavern markets that build, in which Germany's 88~GW
    at 48 full-load hours dominates the weighting while Denmark's 40~per~cent recovery
    carries only 8~GW. Panel~(b) is the cumulative profit of each fleet to the year
    shown, and the shaded wedge between them is what hydrogen loses over and above gas.}
    \includegraphics[width=0.96\textwidth]{paper/P32_peaker_recovery_bridge.pdf}
    \label{fig:peakerbridge}
    \\[2pt]{\footnotesize\textit{Source: Authors' contribution.}}
\end{figure}

\subsection{Hydrogen Supply and Network Infrastructure}
\label{sec:power:infra}

Two infrastructure questions bound the picture without changing it. On \emph{supply}, the
delivered hydrogen cost that enters Equation~\eqref{eq:srmc} is dominated by domestic green
hydrogen. Across the 29 countries the cheapest 2050 route is domestic electrolysis in 23 and
pipeline import in six, at a delivered cost of about \euro{}37 to \euro{}64/MWh (about \euro{}1.24 to
\euro{}2.14/kg), with shipped ammonia and blue or pink hydrogen more expensive
\citep{Buhler2025hydrogenEU,OIES2022hydrogenimports}. A hydrogen economy for
peaking needs only electrolysers, renewable supply and storage, and so avoids the critical-mineral
and lithium dependence of battery backup; its principal cost is the seasonal storage that
Equation~\eqref{eq:srmc} already prices, which we treat as a cited trajectory rather than a built network, and leave, together with ammonia and other carriers, to further work.

On \emph{networks}, each pathway carries a distinct bill (cumulative 2025 to 2050,
central estimate, all 29 markets; Appendix~\ref{sec:results:infrabill} and
Figure~\ref{fig:infrabill}): about \euro{}249~bn under Current Policies, \euro{}325~bn under Stated
Policies, \euro{}506~bn under Net Zero, and \euro{}412~bn under H2 Push. In every pathway the
bill is dominated by district-heat expansion and grid reinforcement, which a hydrogen network does not drive. H2 Push is the cautionary case. It pays about \euro{}74~bn for a hydrogen
distribution network \emph{and} still about \euro{}89~bn of grid reinforcement, because
hydrogen
serves only the peak while heat pumps continue to set the distribution-level load. A hydrogen role
genuinely confined to firm peaking, as the merit order implies, argues for reusing the existing gas
transmission backbone and salt caverns for seasonal balancing \citep{GasForClimate2020backbone,Caglayan2020caverns}, and not for building a parallel hydrogen distribution
grid to dwellings. This is the buildings-to-power synergy in its sharpest form. Whatever role hydrogen holds sits at the firm-capacity end of the power system, and Section~\ref{sec:power:caveat} has already shown that the plant holding it recovers about 9~per~cent of its capital from energy-only rents, so the role depends on a standing payment the market does not currently provide. The synergy is therefore a claim about where hydrogen belongs if it is built, and not a claim that building it pays.
